\section{What can be done and what is open}

\subsection{What can be done}\label{dc:sec:answer}

Partly, and the parts can be stated exactly: the deconvolution of the Thorin density always exists, as Weil's distribution computed from the primes with no hypothesis (Theorem~\ref{dc:thm:weil}), and its positivity to depth $\varepsilon$ is a zero free half plane (Theorem~\ref{dc:thm:damped}), known for $\varepsilon=\frac12$ and announced for $\varepsilon=\frac38$, each rung $\lambda_{\min}(\frac12+\varepsilon,a)\ge0$ being a finite problem in the primes up to $e^{2a}$ whose limit, under the hypotheses of Theorem~\ref{dc:thm:ladder}, is a margin $2b_{1/2+\varepsilon}$ of about $0.092\,\varepsilon$, approached at a rate of order $a^{-2}$.

A rung is finite information in the sense of Section~\ref{sec:ch13:visibility} (Theorem~\ref{thm:ch13:perturb} is the rigorous form of this for the heat trace), so heuristically no finite set of rungs sees every height, and at $\varepsilon=0$ the margin is $0$ and, under RH, the approach to it is faster than any power of $a$ (Theorem~\ref{dc:thm:ladder}(c)): this is the deconvolution obstruction, consistent with the theorem of Rodgers and Tao that RH, if true, holds with no slack~\cite{RodgersTao}.

\subsection{Open problems}\label{dc:sec:open}

\begin{problem}[Certified rungs]\label{dc:prob:cert}
Certify $\lambda_{\min}(\frac12+\varepsilon,a)>0$ in interval arithmetic for $\varepsilon\in\{0.05,0.1,0.2,0.375\}$ and $a=4$, which by monotonicity in $a$ covers every $a\le4$, and determine how far in $a$ a certificate can be pushed at fixed $\varepsilon$. Each certificate is unconditional, and at $\varepsilon=\frac38$ each agrees with what the announced half plane gives, namely $\lambda_{\min}\ge0$ (Theorem~\ref{dc:thm:damped}), and can be checked independently of it.
\end{problem}

\begin{problem}[The exact rate]\label{dc:prob:rate}
Prove that, under the hypotheses of Theorem~\ref{dc:thm:ladder}, $a^2(\lambda_{\min}(\frac12+\varepsilon,a)-2b_{1/2+\varepsilon})\to\pi^2v_{1/2+\varepsilon}''(0)/4$, and find the next term. The converged values of Table~\ref{dc:tab:rate} exceed the constant by a relative amount that falls like $a^{-1}$: a least squares fit on $2\le a\le7$ gives $a^2(\lambda_{\min}(\frac12+\varepsilon,a)-2b_{1/2+\varepsilon})=\bar c_\varepsilon(C_0+C_1/a+C_2/a^2)$, with $\bar c_\varepsilon=\pi^2v_{1/2+\varepsilon}''(0)/4$, $C_0=1.0002$, $1.0003$, $1.0005$, $C_1=0.139$, $0.159$, $0.166$ and $C_2=0.037$, $0.045$, $0.051$ at $\varepsilon=\frac12$, $0.2$, $0.1$, and residuals below $1.1\times10^{-4}\,\bar c_\varepsilon$. So, numerically, the correction is positive and $\bar c_\varepsilon$ is not a bound at finite $a$ (the scaled excess is $1.024\,\bar c_\varepsilon$ at $a=6$, $\varepsilon=\frac12$); the fit is numerical evidence for the limit, to within $0.04$ to $0.2$ percent, and for the order $a^{-1}$ of the correction, and neither is proved.
\end{problem}

\begin{problem}[The transition at the line]\label{dc:prob:transition}
Describe $\lambda_{\min}(\frac12+\varepsilon,a)$ jointly as $a\to\infty$ and $\varepsilon\to0$. For fixed $\varepsilon>0$, under the hypotheses of Theorem~\ref{dc:thm:ladder}, the limit is $2b_{1/2+\varepsilon}$, which is at most $2V_\xi\,\varepsilon$ and asymptotic to it as $\varepsilon\downarrow0$, and the rate is algebraic; at $\varepsilon=0$, under RH, the decay is faster than any power. On which curves $\varepsilon=\varepsilon(a)\to0$ does $\lambda_{\min}(\frac12+\varepsilon,a)$ stay comparable to $2b_{1/2+\varepsilon}$, and does the crossover occur at $\varepsilon\approx$ the mean spacing of the zeros at height $e^{2a}$?
\end{problem}

\begin{problem}[Uniformity from structure]\label{dc:prob:uniform}
Find a structural reason, valid for all $a$, why $Q_\varepsilon\ge0$ on $L^2[-a,a]$ for some fixed $\varepsilon<\frac12$, using properties of the archimedean and pole terms that dominate the primes up to $e^{2a}$ uniformly. By Theorem~\ref{dc:thm:damped} any such argument proves a zero free half plane $\re s>\frac12+\varepsilon$.
\end{problem}
